\section{A Budget-Matched Framework for Test-Time Ensembling}
\label{sec:method}

\subsection{Setup, notation, and suite roles}

\emph{In one sentence:} two question suites in fixed roles --- a development suite on which conditions, strata, and rules were designed, and a held-out suite on which the primary evaluation is conducted; all comparisons hold the question, the model, and the realized decoding randomness constant by construction.

Let $M$ denote a large language model and let $q$ denote a question with a gold answer $y^*(q)$. For multiple-choice tasks the answer is an option letter, and for free-form numeric tasks the answer is a value compared numerically. A prompt $x$ is a rendering of $q$ into natural language, consisting of an optional passage, a question stem, an option list, and an instruction that specifies the required output format. We distinguish one canonical prompt $x_0$ and $K$ paraphrased prompts $x_1, \dots, x_K$ that preserve the semantics of $x_0$ under a validation procedure described in Section~\ref{sec:setup}. The study uses two question suites in fixed roles: a \emph{development suite} of 90 questions on which conditions, strata, and allocation rules were designed, and a \emph{held-out suite} of 119 fresh questions, drawn mechanically from the same pools by replaying the seeded shuffle (Section~\ref{sec:setup}) without consulting any model output, on which the primary evaluation is conducted. Because the strata and rules were fixed after seeing the development-suite results, the held-out suite is a transfer test of previously chosen designs, not a preregistered confirmation; we use the term \emph{held-out} rather than blind throughout.

The model produces an answer through stochastic decoding. We write $a(M, x, s)$ for the answer extracted from the $s$-th independent decode of prompt $x$ at temperature $\tau = 0.7$, and $a_g(M, x)$ for the greedy decode at $\tau = 0$. Answers are extracted from the model output by a fixed parser that searches for the mandated answer format and falls back to pattern matching over the final output lines; the parse rate is reported in Section~\ref{sec:setup}, together with a retry policy and a sensitivity analysis for failed extractions.

\subsection{The call matrix}

\emph{In one sentence:} every model--question pair yields one sixteen-call matrix from which all experimental conditions are derived deterministically.

For every question and model we collect a matrix of sixteen calls:
\begin{equation}
\mathcal{C}(M, q) \;=\; \bigl\{ a_g(M, x_0) \bigr\} \;\cup\; \bigcup_{s=1}^{S} \bigl\{ a(M, x_0, s) \bigr\} \;\cup\; \bigcup_{k=1}^{K} \bigl\{ a(M, x_k, s) : s \le r_k \bigr\},
\label{eq:matrix}
\end{equation}
with $S = 6$ samples on the canonical prompt, $K = 6$ paraphrases, and $r_k = 2$ decodes for the first three paraphrases and $r_k = 1$ for the remaining three. The matrix therefore contains one greedy call, six canonical samples, six first-decode paraphrase calls, and three second-decode paraphrase calls. Every experimental condition studied in this paper is a deterministic function of this matrix, which guarantees that all comparisons hold the question, the model, and the realized decoding randomness constant across conditions.

\subsection{Ensemble conditions}
\label{sec:conditions}

\emph{In one sentence:} all conditions aggregate votes by majority rule, both six-vote conditions spend exactly six calls at the same decoding temperature, and the paraphrase ensemble mixes prompt variation with per-call sampling noise --- the latter is stated wherever paraphrase margins are interpreted.

All conditions aggregate votes by majority rule. Ties are broken in favor of the earliest vote in a fixed canonical order, which prefers the canonical prompt over paraphrases and earlier paraphrases over later ones; we report tie incidence, which is negligible except where noted.

The \emph{self-consistency} condition of size $k$ aggregates the first $k$ canonical samples:
\begin{equation}
\mathrm{SC}\text{-}k(q) \;=\; \operatorname*{majority}\bigl( a(M, x_0, 1), \dots, a(M, x_0, k) \bigr), \qquad k \le 6 .
\end{equation}
The \emph{paraphrase ensemble} of size $k$ aggregates one decode of each of the first $k$ paraphrases:
\begin{equation}
\mathrm{PE}\text{-}k(q) \;=\; \operatorname*{majority}\bigl( a(M, x_1, 1), \dots, a(M, x_k, 1) \bigr), \qquad k \le 6 .
\end{equation}
Both conditions spend exactly $k$ model calls at the same decoding temperature ($\tau = 0.7$), so every vote in either condition is an independent stochastic draw. They differ in which factor is free to vary: in $\mathrm{SC}\text{-}k$ the prompt is fixed and only the decoding draw varies, whereas in $\mathrm{PE}\text{-}k$ the prompt varies and \emph{each paraphrase call additionally carries the same per-call decoding noise as a canonical sample, on top of its prompt variation}. A PE-$k$ margin over an SC-$k$ margin therefore characterizes paraphrase ensembles \emph{with sampling noise mixed in}, not pure prompt-space effects; isolating the prompt component requires re-decoding each paraphrase greedily, a control we collected on the development suite only (\ref{app:greedypara}). Wherever this paper interprets a paraphrase-ensemble margin, the statement is qualified accordingly. The budget match is exact in calls and identical in decoding configuration --- this is not a comparison of stochastic against deterministic decoding.

The \emph{hybrid} condition splits a budget of six calls between the two sources, using two decodes of each of the first three paraphrases:
\begin{equation}
\mathrm{HB}\text{-}6(q) \;=\; \operatorname*{majority}\bigl( a(M, x_k, s) : k \in \{1,2,3\},\ s \in \{1,2\} \bigr).
\end{equation}
Finally, the \emph{combined} condition aggregates all twelve stochastic votes available in the matrix, $\mathrm{COMB}\text{-}12(q) = \operatorname{majority}\bigl( \mathrm{SC}\text{-}6(q) \cup \mathrm{PE}\text{-}6(q) \bigr)$, and serves as a twelve-call reference point for the marginal value of doubling the budget.

Three baselines and controls complete the design. The \emph{single-sample} baseline $\mathrm{S1}$ is one stochastic draw of the canonical prompt at the decoding temperature used throughout ($\tau = 0.7$); its expected accuracy is estimated by averaging over the six canonical samples, and it is the \emph{primary} $k{=}1$ baseline of this study, because it is the marginal call that the ensembles actually add to. Its paraphrase-side analogue, $\mathrm{PE}\text{-}1$, is the expected accuracy of one paraphrase draw, estimated by averaging over the six first-decode paraphrase calls; the difference $\mathrm{PE}\text{-}1 - \mathrm{S1}$ measures the \emph{paraphrase difficulty margin} --- how much easier the paraphrased prompts are at a single call --- and the within-condition aggregation gains are $\Delta\mathrm{SC} = \mathrm{SC}\text{-}6 - \mathrm{S1}$ and $\Delta\mathrm{PE} = \mathrm{PE}\text{-}6 - \mathrm{PE}\text{-}1$. All gain statements below state explicitly which of these references they use. The greedy decode $a_g(M, x_0)$ at $\tau = 0$ is reported as a \emph{reference} baseline: it is a different decoding configuration, not a draw from the same distribution, and \ref{app:dev} quantifies how far it can sit from S1. A \emph{greedy-paraphrase} control re-decodes each paraphrase at $\tau = 0$, giving $\mathrm{PE}^{g}\text{-}6(q) = \operatorname{majority}\bigl( a_g(M, x_1), \dots, a_g(M, x_6) \bigr)$; comparing $\mathrm{PE}^{g}\text{-}6$ with $\mathrm{PE}\text{-}6$ isolates the contribution of paraphrase variation from the decoding stochasticity that the standard paraphrase ensemble mixes in. The control was collected on the development suite only (Section~\ref{sec:setup}), so held-out paraphrase margins retain that mixing and are read with the caveat stated where they appear. A chain-of-thought variant repeats the single-call baseline, the self-consistency condition, and (for stem-locus datasets) the paraphrase ensemble under an explicit reasoning instruction, and a hard-task probe repeats the single-call and self-consistency conditions on MMLU-Pro; both are specified in Section~\ref{sec:setup} and analyzed in \ref{app:greedypara}--\ref{app:cot}.

\subsection{Consistency, modal share, and fragility}

\emph{In one sentence:} three per-question statistics --- sample agreement, paraphrase agreement, and fragility --- summarize the two answer distributions and their disagreement.

\emph{Sample agreement} is the frequency of the modal answer among the six canonical samples,
\begin{equation}
c_{\mathrm{S}}(q) \;=\; \max_{v} \frac{1}{6} \sum_{s=1}^{6} \mathbb{1}\bigl[ a(M, x_0, s) = v \bigr],
\end{equation}
and \emph{paraphrase agreement} $c_{\mathrm{P}}(q)$ is defined analogously over the six first-decode paraphrase answers. Both lie in $(0, 1]$, and the value $1$ indicates that the corresponding source produced no diversity at all, a situation we call \emph{answer degeneracy}. The \emph{modal vote share} is the corresponding statistic over the union of the twelve stochastic votes, that is, the fraction of all twelve votes that fall on the modal answer, and it is our primary per-question disagreement measure: it reads directly as the strength of the majority. We also compute the support-normalized answer entropy of the same twelve-vote distribution and report it alongside the modal share for the follow-up conditions, where normalized entropy is the conventional statistic; all entropy values in the paper use this normalization and are reported as fractions, not percentages. The modal share is preferred as the headline statistic because a support-normalized entropy inflates near-unanimous distributions: with twelve votes, a single dissenter already yields a normalized entropy of $0.42$, so any entropy threshold below that value merely separates degenerate from non-degenerate questions rather than grading the strength of the majority.

\emph{Fragility} captures disagreement between the two views of the same question: a question is fragile for a model when the majority answer of the canonical samples differs from the majority answer of the paraphrase answers, that is, when $\mathrm{SC}\text{-}6(q) \ne \mathrm{PE}\text{-}6(q)$. Fragility measures how much of the model's answer is an artifact of the particular surface form of the prompt rather than of the underlying question. The greedy-paraphrase control supports a noise-free variant, $\mathrm{SC}\text{-}6(q) \ne \mathrm{PE}^{g}\text{-}6(q)$, in which neither side involves stochastic decoding.

\subsection{Error detection and selective prediction}

\emph{In one sentence:} agreement scores are evaluated as unsupervised error detectors by AUROC and tie-randomized AURC, with the degenerate stratum reported separately because detector outcomes there are definitional.

For each model, the consistency score of a question (reversed, so that lower agreement indicates higher predicted risk) ranks questions, and we compute the area under the receiver operating characteristic curve (AUROC) for predicting whether the reference answer is incorrect, where the reference is the combined twelve-vote ensemble unless stated otherwise. When both agreement scores are available, the \emph{combined} detector is their sum of disagreements, $2 - c_{\mathrm{S}} - c_{\mathrm{P}}$, so that a question is flagged when either source shows diversity. We complement AUROC with risk-coverage curves summarized by the area under the risk-coverage curve (AURC), computed by accepting questions in decreasing order of confidence and averaging the error rate over all prefix sizes; smaller AURC indicates a better ranking, and a random ranking yields AURC equal to the overall error rate. Because several questions typically share the same confidence value, we average AURC over uniformly random tie-breaking (1{,}000 permutations). Where a model produces zero errors or constant scores, AUROC is undefined and we report the regime instead; where a model produces fewer than five errors, we report AUROC and AURC descriptively only, without confidence intervals or cross-model comparison --- a rule that proved necessary: a development-suite reading that looked decisive at one to three errors did not survive the held-out recomputation at twelve to twenty-one errors (Section~\ref{sec:heldout}).

Because the agreement scores and the reference answer share the same underlying votes, part of the error-detection result is structural rather than empirical: when all twelve votes coincide, the reference equals that unanimous answer by construction, and any detector that reads unanimity is correct by definition on such questions. We therefore separate the degenerate stratum, where the detector outcome is definitional, from the non-degenerate stratum, where the ranking carries empirical content, and we report both. As an additional check against circularity, we recompute the entropy-gain correlation using disjoint votes.

\subsection{Statistical analysis}
\label{sec:stats}

\emph{In one sentence:} accuracy comparisons are paired by construction and multiplicity is handled by designation --- four designated primary contrasts on the held-out suite, everything else exploratory.

Accuracy comparisons between two conditions on the same questions are paired by construction. We report bootstrap 95\% confidence intervals for accuracy differences (10,000 resamples, resampling questions) and two-sided exact McNemar tests for paired correctness indicators. Because we interpret a difference of three percentage points or less as practically negligible, we also report two one-sided tests (TOST) against equivalence bounds of $\pm 3$ points, implemented as bootstrap 90\% confidence intervals for the paired difference: equivalence is established only when the entire interval lies inside the bounds. Multiplicity is handled by designation rather than by blanket correction: the four contrasts on the held-out suite that the study is designed to answer are designated \emph{primary} (the pooled gains of SC-6, PE-6, and COMB-12 over S1, and the PE-6 versus SC-6 gap); their two-sided bootstrap $p$-values are Holm-adjusted within the family and reported alongside the confidence intervals. Every other comparison in the paper --- per-model contrasts, per-dataset cells, stratum decompositions, and the entire development suite --- is labeled exploratory and carries no correction claim. One caveat is inherent to the design: accuracy moves in units of one question, so the $\pm 3$-point bound corresponds to 2.7 questions at $n = 90$ and fewer than one question for the smallest completed sets, which makes the TOST verdict on the small models sensitive to the flip of a single question; we flag such cells explicitly. A failure to reject the null hypothesis of no difference is reported as \emph{no detectable difference} and never as evidence of equivalence; at our sample sizes, differences smaller than roughly five points often remain undetectable, and we state this power limitation explicitly wherever it applies. Pooled analyses aggregate (model, question) pairs that share questions across models; because such pairs are not independent, all pooled confidence intervals use a question-level cluster bootstrap that resamples questions and carries every model's pair within a resampled question. Per-model analyses use each model's own completed question set, with the sample size stated in every table. We treat question sets as fixed and do not extrapolate beyond the measured benchmarks.

\subsection{When does majority voting help? A heuristic model}
\label{sec:condorcet}

\emph{In one sentence:} a plurality-vote model in the spirit of the Condorcet jury theorem \citep{condorcet1785essai} explains the regime structure --- voting pays where the modal answer is correct, amplifies errors where it is wrong --- and justifies taking S1, not greedy, as the $k{=}1$ baseline.

We state it as a heuristic, not a theorem: the asymptotic statements are standard, but the finite-$k$ behavior, the tie handling, and the application to language-model votes each require idealizations that we do not prove. Write $p$ for the probability that one vote drawn from the canonical sampling distribution is the correct answer, and $q_1 \ge q_2 \ge \dots$ for the probabilities of the incorrect answers, so that $p + \sum_i q_i = 1$. Three idealizations are implicit: the $k$ votes are independent and identically distributed draws from this distribution; plurality ties are broken uniformly at random; and the paraphrase ensemble is described by an analogous distribution of its own, which shares the same gold answer but need not coincide with the canonical one, because paraphrased prompts are not exchangeable with resampled decodes of the original prompt.

\textbf{Asymptotic statement.} As $k$ grows, the plurality winner converges in probability to the modal answer of the vote distribution. Consequently, when $p > q_1$ the probability that the plurality winner is correct tends to one, and when $q_1 > p$ it tends to zero. The convergence is asymptotic; we do not assert monotonicity in $k$, and at small $k$ with near-tied leaders the win probability can move in either direction, which is visible in our budget curves as one-question zigzags (\ref{app:dev}).

\textbf{Baseline statement.} A single vote is correct with probability $p$, so majority voting with $k \to \infty$ outperforms \emph{a single stochastic sample} if and only if $p > q_1$, and it underperforms it exactly when the modal answer is wrong. The comparison against greedy decoding is subtler and not covered by the model: greedy decoding approximates the argmax of the model's answer distribution rather than drawing one sample from it, so the greedy answer and the modal sampled answer coincide only when temperature-$0$ decoding is faithful to the underlying distribution. \ref{app:dev} shows that the two can disagree by up to ten points on real deployments, which is precisely a violation of this idealization, and it is why we take the single-sample expectation S1, not greedy, as the $k{=}1$ baseline.

The model explains the three regimes we observe. In the degenerate regime, $p$ is zero or one and voting is vacuous. In the high-entropy regime with $p > q_1$, self-consistency converts a distribution that is right on average into a reliable majority. In the high-entropy regime with $q_1 > p$, voting amplifies the modal error and reduces accuracy relative to a single sample, while a second diversity source that shifts the effective vote distribution can restore $p > q_1$. The model also clarifies why disagreement alone does not predict gains: the modal share measures the spread of the distribution, whereas the sign of $p - q_1$ measures whether the mode is correct, and two distributions with identical shares can sit on opposite sides of that threshold. The budget-matched comparison of Sections~\ref{sec:heldout} and \ref{app:dev} can then be read as asking which diversity source moves the effective vote distribution further toward the $p > q_1$ region at a given budget.
